% !PRINTED 575
\weberpage{575}
%<<CONT
consideration in (3) $\mathfrak{m}$ can denote any arbitrary ideal in $\Omega$, we arrive at the following \gesperrt{principal theorem}:

\begin{quote}
% [sic] The source says in theorem 4 that T_m is in Omega; theorem 3 establishes coefficient membership only in the class field. This unsupported narrowing of the base field is retained as a source discrepancy.
4. \gesperrt{If $\mathfrak{m}$ is an arbitrary ideal of the field $\Omega$, then there exists a function $T_{\mathfrak m}$ in $\Omega$ whose roots are those of the numbers $\tau_\xi$ that are distinct from one another, when $\xi$ runs through a system of incongruent numbers coprime to $\mathfrak{m}$.}
\end{quote}

To determine the degree of the function $T_{\mathfrak m}$, we observe that two quantities $\tau_\xi, \tau_{\xi'}$ are equal to one another only when
\weq{8}{\xi \equiv \varepsilon\xi' \quad (\mathrm{mod}\ \mathfrak{m})}
holds. Thus just as many of the $\tau_\xi$ become equal to one another as there are units incongruent modulo $\mathfrak{m}$; in general these are two, for $\Delta = -4$ they are four, and for $\Delta = -3$ they are six. These numbers are, however, reduced still further when different ones among the units $\varepsilon$ are congruent with respect to the modulus $\mathfrak{m}$, that is, when $1 - \varepsilon$ is divisible by $\mathfrak{m}$. But this can occur only when $\mathfrak{m}$ is a divisor of 2 or of 3. Hence:

\begin{quote}
5. \gesperrt{The degree of the function $T_{\mathfrak m}$ is equal to the number $\psi(\mathfrak{m})$} (Bd. II, \S\ 168) \gesperrt{of the numbers in $\Omega$ that are incongruent with respect to the modulus $\mathfrak{m}$ and coprime to $\mathfrak{m}$, divided by the number of units in $\Omega$ that are incongruent with respect to $\mathfrak{m}$.}
\end{quote}

The roots of the function $T_{\mathfrak m}$ determine an algebraic field $\mathfrak{T}_{\mathfrak m}$ over the class field, whose relative degree is at most equal to the degree of $T_{\mathfrak m}$, and both degrees are equal when $T_{\mathfrak m}$ is irreducible.

\begin{quote}
6. \gesperrt{This field $\mathfrak{T}_{\mathfrak m}$ we shall call the division field for the modulus $\mathfrak{m}$.}
\end{quote}

These division fields play a similar role for the field $\Omega$ to that played by the cyclotomic fields for the field of rational numbers, except that here the class field is still interposed between them, and to this there is no analogue in the field of rational numbers, just as there is none in any one-class field.

By the multiplication theorem (\S\ 151),
\weq{9}{\tau(\xi u) = \theta_\xi[\tau(u)]}
can be expressed rationally in terms of $\tau(u)$. If we set $u = \xi'\dfrac{\alpha}{\mu}$, then it follows:
\weq{10}{\tau_{\xi\xi'} = \tau\left(\frac{\xi\xi'\alpha}{\mu}\right) = \theta_\xi(\tau_{\xi'}),}
%CONT>>
